\documentclass{beamer}

\pdfmapfile{+sansmathaccent.map}


\mode<presentation>
{
	\usetheme{Warsaw} % or try Darmstadt, Madrid, Warsaw, Rochester, CambridgeUS, ...
	\usecolortheme{seahorse} % or try seahorse, beaver, crane, wolverine, ...
	\usefonttheme{serif}  % or try serif, structurebold, ...
	\setbeamertemplate{navigation symbols}{}
	\setbeamertemplate{caption}[numbered]
} 


%%%%%%%%%%%%%%%%%%%%%%%%%%%%
% itemize settings


%%%%%%%%%%%%%%%%%%%%%%%%%%%%
% itemize settings

\definecolor{myhotpink}{RGB}{255, 80, 200}
\definecolor{mywarmpink}{RGB}{255, 60, 160}
\definecolor{mylightpink}{RGB}{255, 80, 200}
\definecolor{mypink}{RGB}{255, 30, 80}
\definecolor{mydarkpink}{RGB}{155, 25, 60}

\definecolor{mypaleblue}{RGB}{240, 240, 255}
\definecolor{mylightblue}{RGB}{120, 150, 255}
\definecolor{myblue}{RGB}{90, 90, 255}
\definecolor{mygblue}{RGB}{70, 110, 240}
\definecolor{mycourseblue}{RGB}{120, 140, 255}
\definecolor{mycourselightblue}{RGB}{245, 250, 255}
\definecolor{mydarkblue}{RGB}{0, 0, 180}
\definecolor{myblackblue}{RGB}{40, 40, 120}

\definecolor{mygreen}{RGB}{0, 200, 0}
\definecolor{mygreen2}{RGB}{245, 255, 230}
\definecolor{mydarkgreen}{RGB}{0, 120, 0}
\definecolor{mydarkgreen}{RGB}{0, 120, 0}


\definecolor{mygray}{gray}{0.8}
\definecolor{mydarkgray}{RGB}{80, 80, 160}
\definecolor{mylightgray}{RGB}{160, 160, 160}

\definecolor{mydarkred}{RGB}{160, 30, 30}
\definecolor{mylightred}{RGB}{255, 150, 150}
\definecolor{myred}{RGB}{200, 110, 110}
\definecolor{myblackred}{RGB}{120, 40, 40}

\definecolor{mypink}{RGB}{255, 30, 80}
\definecolor{myhotpink}{RGB}{255, 80, 200}
\definecolor{mywarmpink}{RGB}{255, 60, 160}
\definecolor{mylightpink}{RGB}{255, 80, 200}
\definecolor{mydarkpink}{RGB}{155, 25, 60}
\definecolor{mywhitepink}{RGB}{255, 240, 240}

\definecolor{mydarkcolor}{RGB}{60, 25, 155}
\definecolor{mylightcolor}{RGB}{130, 180, 250}



\setbeamertemplate{itemize items}[default]

\setbeamertemplate{itemize item}{\color{myblackblue}$\blacksquare$}
\setbeamertemplate{itemize subitem}{\color{mygblue}$\blacktriangleright$}
\setbeamertemplate{itemize subsubitem}{\color{mygray}$\blacksquare$}

\setbeamercolor{palette quaternary}{fg=white,bg=mycourseblue}
\setbeamercolor{titlelike}{parent=palette quaternary}

\setbeamercolor{palette quaternary2}{fg=black,bg=mycourselightblue}
\setbeamercolor{frametitle}{parent=palette quaternary2}

\setbeamerfont{frametitle}{size=\Large,series=\scshape}
\setbeamerfont{framesubtitle}{size=\normalsize,series=\upshape}





%%%%%%%%%%%%%%%%%%%%%%%%%%%%
% block settings

\setbeamercolor{block title}{bg=red!30,fg=black}

\setbeamercolor*{block title example}{bg=mygreen!40!white,fg=black}

\setbeamercolor*{block body example}{fg= black, bg= mygreen2}


%%%%%%%%%%%%%%%%%%%%%%%%%%%%
% URL settings
\hypersetup{
	colorlinks=true,
	linkcolor=blue,
	filecolor=blue,      
	urlcolor=blue,
}

%%%%%%%%%%%%%%%%%%%%%%%%%%

\renewcommand{\familydefault}{\rmdefault}


\newcommand{\mycolor}[1] {
\setbeamercolor{normal text}{fg={#1}} % Recolors standard body text
\setbeamercolor{structure}{fg={#1}}   % Recolors bullets and structure items
\usebeamercolor[fg]{normal text}      % Applies the normal text color immediately
}


\usepackage{amsmath}
\usepackage{mathtools}

\usepackage{cancel}
\usepackage{centernot}

\usepackage{changepage}
\usepackage{subcaption}

\usepackage{qrcode}

\newcommand{\bo}[1] {\mathbf{#1}}
\newcommand{\R}{\mathbb{R}} 
\newcommand{\T}{^\top}     



\newcommand{\mydate}{Fall 2026}

\newcommand{\mygit}{\textcolor{blue}{\href{https://github.com/SergeiSa/Advanced-Control-2026}{github.com/SergeiSa/Advanced-Control-2026}}}

\newcommand{\myqr}{ \textcolor{black}{\qrcode[height=1.5in]{https://github.com/SergeiSa/Advanced-Control-2026}}
}

\newcommand{\myqrframe}{
	\begin{frame}
		\centerline{Lecture slides are available via Github:}
		\bigskip
		\centerline{\mygit}
		\bigskip
		\myqr
	\end{frame}
}


\newcommand{\bref}[2] {\textcolor{blue}{\href{#1}{#2}}}

%%%%%%%%%%%%%%%%%%%%%%%%%%%%
% code settings

\usepackage{listings}
\usepackage{color}
% \definecolor{mygreen}{rgb}{0,0.6,0}
% \definecolor{mygray}{rgb}{0.5,0.5,0.5}
\definecolor{mymauve}{rgb}{0.58,0,0.82}
\lstset{ 
	backgroundcolor=\color{white},   % choose the background color; you must add \usepackage{color} or \usepackage{xcolor}; should come as last argument
	basicstyle=\footnotesize,        % the size of the fonts that are used for the code
	breakatwhitespace=false,         % sets if automatic breaks should only happen at whitespace
	breaklines=true,                 % sets automatic line breaking
	captionpos=b,                    % sets the caption-position to bottom
	commentstyle=\color{mygreen},    % comment style
	deletekeywords={...},            % if you want to delete keywords from the given language
	escapeinside={\%*}{*)},          % if you want to add LaTeX within your code
	extendedchars=true,              % lets you use non-ASCII characters; for 8-bits encodings only, does not work with UTF-8
	firstnumber=0000,                % start line enumeration with line 0000
	frame=single,	                   % adds a frame around the code
	keepspaces=true,                 % keeps spaces in text, useful for keeping indentation of code (possibly needs columns=flexible)
	keywordstyle=\color{blue},       % keyword style
	language=Octave,                 % the language of the code
	morekeywords={*,...},            % if you want to add more keywords to the set
	numbers=left,                    % where to put the line-numbers; possible values are (none, left, right)
	numbersep=5pt,                   % how far the line-numbers are from the code
	numberstyle=\tiny\color{mygray}, % the style that is used for the line-numbers
	rulecolor=\color{black},         % if not set, the frame-color may be changed on line-breaks within not-black text (e.g. comments (green here))
	showspaces=false,                % show spaces everywhere adding particular underscores; it overrides 'showstringspaces'
	showstringspaces=false,          % underline spaces within strings only
	showtabs=false,                  % show tabs within strings adding particular underscores
	stepnumber=2,                    % the step between two line-numbers. If it's 1, each line will be numbered
	stringstyle=\color{mymauve},     % string literal style
	tabsize=2,	                   % sets default tabsize to 2 spaces
	title=\lstname                   % show the filename of files included with \lstinputlisting; also try caption instead of title
}


%%%%%%%%%%%%%%%%%%%%%%%%%%%%
% URL settings
\hypersetup{
	colorlinks=false,
	linkcolor=blue,
	filecolor=blue,      
	urlcolor=blue,
}

%%%%%%%%%%%%%%%%%%%%%%%%%%

%%%%%%%%%%%%%%%%%%%%%%%%%%%%
% tikz settings

\usepackage{tikz}
\tikzset{every picture/.style={line width=0.75pt}}
\usetikzlibrary{shapes.geometric, decorations.markings, positioning, patterns, calc, arrows.meta}


\subtitle{Advanced Control}
\author{by Sergei Savin}
\centering
\date{\mydate}



\title{Stochastic Bellman Equation}
\begin{document}
\maketitle


\begin{frame}{Content}
	\begin{itemize}
		\item Stochastic sequences and expectation
		\item Stochastic linear systems
		\item Finite-horizon stochastic cost
		\item Stochastic Bellman equation
		\item Stochastic Riccati equation
		\item Infinite-horizon cost with discounting
		\item Value functions and action-value functions
		\item Discounted Bellman and Riccati equations
		\item Appendix: derivation of the stochastic Riccati equation
	\end{itemize}
\end{frame}



%------------------------------------------------
\begin{frame}{Stochastic sequences}
	\begin{flushleft}
		

Let $w_k$ be a random (stochastic) sequence with expectation
\[
\mu = \mathbb{E}[w_k].
\]

Given $N$ samples, we can estimate its expectation using the sample mean:
\[
\hat{\mu}_N = \frac{1}{N}\sum_{i=k}^{k+N-1} w_i.
\]

Under appropriate assumptions (e.g., independent identically distributed
samples with finite mean), we have
\[
\lim_{N\to\infty}\hat{\mu}_N = \mathbb{E}[w_k].
\]

\vspace{0.3cm}

Expectation is a linear operator. In particular, for constant matrices
$A,B$ and random vectors $x,u$:
\[
\boxed{
	\mathbb{E}[Ax+Bu]
	= A\mathbb{E}[x]+B\mathbb{E}[u]
}
\]
		
	\end{flushleft}
\end{frame}
%------------------------------------------------

\begin{frame}{Stochastic linear systems}
	\begin{flushleft}
		
Consider a discrete-time LTI system:
\[
x_{k+1}=Ax_k+w_k,
\]
where $w_k$ is an element of a stochastic sequence.

We can make the following observations:

\begin{itemize}
	\item Without knowing the realization of $w_k$, we cannot predict
	the exact value of $x_{k+1}$.
	
	\item Consequently, even for a known initial state, $x_{k+1}$
	is a random variable.
	
	\item While exact state values cannot be predicted, we can
	analyze their statistical properties.
\end{itemize}

\vspace{0.2cm}

For example, using the linearity of expectation:
\[
\boxed{
	\mathbb{E}[x_{k+1}]
	= A\mathbb{E}[x_k]+\mathbb{E}[w_k]
}
\]
		
	\end{flushleft}
\end{frame}
%------------------------------------------------

\begin{frame}{Finite-horizon stochastic cost}
	\begin{flushleft}
		
Let us introduce a quadratic cost function suitable for a
stochastic linear system:
\[
J = \mathbb{E}\left[
\sum_{k=0}^{N-1} \ell(x_k,u_k)
+ x_N^\top Q_N x_N
\right],
\]
where the stage cost is
\[
\ell(x_k,u_k)
= x_k^\top Q x_k + u_k^\top R u_k.
\]

Here, $Q\succeq 0, Q_N\succeq 0$ and $R\succ 0$.

\vspace{0.3cm}

As in the deterministic case, we define the value function
at step $k$ as the minimum expected cost-to-go:
\[
\boxed{
	V_k(x) =
	\min_{u_k,\ldots,u_{N-1}}
	\mathbb{E}\left[
	\sum_{i=k}^{N-1}\ell(x_i,u_i)
	+x_N^\top Q_N x_N
	\,\middle|\, x_k=x
	\right]
}
\]
		
	\end{flushleft}
\end{frame}
%------------------------------------------------




%------------------------------------------------
\begin{frame}{Stochastic Bellman equation}
	\begin{flushleft}
		
As in the deterministic case, we can use Bellman's principle
to express the value function recursively:

\[
V_k(x_k) =
\min_{u_k}
\left\{
\ell(x_k,u_k)
+ \mathbb{E}\left[
V_{k+1}(x_{k+1}) \mid x_k,u_k
\right]
\right\}.
\]

For our stochastic linear system:
\[
x_{k+1}=Ax_k+Bu_k+w_k,
\]
we obtain
\[
\boxed{
	V_k(x) =
	\min_u
	\left\{
	x^\top Qx+u^\top Ru
	+\mathbb{E}\left[
	V_{k+1}(Ax+Bu+w_k)
	\right]
	\right\}.
}
\]

The terminal condition is
\[
V_N(x)=x^\top Q_Nx.
\]
		
	\end{flushleft}
\end{frame}
%------------------------------------------------

\begin{frame}{Stochastic Riccati equation}
	\begin{flushleft}
For zero-mean independent disturbances with covariance
\[
\mathbb{E}[w_k]=0,\qquad
\mathbb{E}[w_kw_k^\top]=W,
\]
the value function retains a quadratic form:
\[
\boxed{
	V_k(x)=x^\top P_kx+c_k.
}
\]

Substituting this expression into the Bellman equation gives
an optimal linear control law:
\[
u_k=K_kx_k,
\]
where
\[
	K_k =
	-(R+B^\top P_{k+1}B)^{-1}
	B^\top P_{k+1}A.
\]

The matrix $P_k$ is found using the Riccati recursion:
\[
		P_k=Q+A^\top P_{k+1}A
		-A^\top P_{k+1}B
		(R+B^\top P_{k+1}B)^{-1}
		B^\top P_{k+1}A.
\]

with terminal condition $P_N=Q_N$.
	\end{flushleft}
\end{frame}
%------------------------------------------------


\begin{frame}{Infinite-horizon cost}
	\begin{flushleft}
		
		{\small 
		We can consider an infinite-horizon version of the stochastic
		optimal control problem.
		
		One issue is that the accumulated cost can become infinite:
		even if the system is stable, persistent stochastic disturbances
		keep introducing uncertainty into the state.
		
		\vspace{0.2cm}
		
		One way to manage this problem is to introduce a discounting
		factor $0 < \gamma < 1$:
		
		\[
		\boxed{
			J = \mathbb{E}\left[
			\sum_{k=0}^{\infty}
			\gamma^k \ell(x_k,u_k)
			\right].
		}
		\]
		
		\vspace{0.2cm}
		
		\begin{itemize}
			\item Discounting allows the expected accumulated cost to remain
			finite under standard bounded-cost assumptions.
			
			\item It gives priority to immediate costs, while progressively
			reducing the importance of costs incurred further into the future.
			
			\item As before, the infinite-horizon formulation has no terminal
			cost component.
		\end{itemize}
	}
	\end{flushleft}
\end{frame}


\begin{frame}{Value function with discounted cost}
	\begin{flushleft}
		
		For a given policy $\pi$, we define the on-policy state-value
		function as the expected discounted cost-to-go:
		
		\[
		V^\pi(x) =
		\mathbb{E}_\pi\left[
		\sum_{k=0}^{\infty}
		\gamma^k \ell(x_k,u_k)
		\,\middle|\, x_0=x
		\right].
		\]
		
		\vspace{0.3cm}
		
		The optimal state-value function is the minimum over all
		admissible policies:
		
		\[
		V^*(x)=\min_\pi V^\pi(x).
		\]
		
	\end{flushleft}
\end{frame}

\begin{frame}{Bellman equation with discounted cost}
	\begin{flushleft}
		
		The infinite-horizon Bellman optimality equation takes the form:
		%
		\[
		V^*(x)=
		\min_u\left\{
		\ell(x,u)+
		\gamma\mathbb{E}\left[
		V^*(x_{k+1})\mid x_k=x,u_k=u
		\right]
		\right\}.
		\]
		%
		We can introduce an action-value function $\mathcal{Q}^\pi(x,u)$,
		which describes the expected cost-to-go when we apply action $u$
		in state $x$, and then follow policy $\pi$:
		%
		\[
		\mathcal{Q}^\pi(x,u)=
		\mathbb{E}_\pi\left[
		\sum_{k=0}^{\infty}\gamma^k\ell(x_k,u_k)
		\,\middle|\,
		x_0=x,\ u_0=u
		\right].
		\]
		%
		The state-value and action-value functions are related by:
		%
		\[
		V^\pi(x)=
		\mathbb{E}_{u\sim\pi(\cdot|x)}
		[\mathcal{Q}^\pi(x,u)].
		\]
		%
		For the optimal policy:
		%
		\[
		V^*(x)=\min_u \mathcal{Q}^*(x,u).
		\]
		
	\end{flushleft}
\end{frame}


\begin{frame}{Riccati equation with discounted cost}
	\begin{flushleft}
		
		For stochastic LQR, the optimal value function is still quadratic:
		\[
		V^*(x)=x^\top Px+c.
		\]
		
		Substituting this expression into the Bellman equation gives
		the discounted algebraic Riccati equation:
		
		\[
				P=Q+\gamma A^\top PA
-\gamma^2 A^\top PB
(R+\gamma B^\top PB)^{-1}
B^\top PA.
		\]
		
		The optimal control law is linear:
		\[
			u^*=-Kx,\qquad
K=(R+\gamma B^\top PB)^{-1}
\gamma B^\top PA.
		\]
		
		The constant term accounts for the stochastic disturbances:
		\[
			c=\frac{\gamma}{1-\gamma}
\operatorname{tr}(PW).
		\]
		
	\end{flushleft}
\end{frame}



\begin{frame}{Literature}
	\begin{itemize}
		
		\item MIT OpenCourseWare:
		\bref{https://ocw.mit.edu/courses/16-323-principles-of-optimal-control-spring-2008/pages/lecture-notes/}
		{Principles of Optimal Control: Lecture 3, Dynamic Programming and Discrete LQR.}
		
		\item Russ Tedrake -- MIT:
		\bref{https://underactuated.mit.edu/lqr.html}
		{Underactuated Robotics: Linear Quadratic Regulators, including the discrete-time Riccati equation.}
		
		\item University of Washington:
		\bref{https://faculty.washington.edu/ratliffl/teaching/2020_LinearSystems.pdf}
		{Lecture 10: Linear Quadratic Regulator -- Discrete Time.}
		
		\item SciPy documentation:
		\bref{https://docs.scipy.org/doc/scipy/reference/generated/scipy.linalg.solve_discrete_are.html}
		{Linear Algebra: \texttt{solve\_discrete\_are}, numerical solution of the discrete-time algebraic Riccati equation.}
		
	\end{itemize}
\end{frame}


\myqrframe

\begin{frame}{Appendix: Derivation of the stochastic Riccati equation}
	\begin{flushleft}
	Substitute the quadratic value function into Bellman's equation:
	\[
	\begin{aligned}
		V_k(x)=\min_u\big\{&
		x^\top Qx+u^\top Ru\\
		&+\mathbb{E}[
		(Ax+Bu+w_k)^\top P_{k+1}
		(Ax+Bu+w_k)]\\
		&+c_{k+1}\big\}.
	\end{aligned}
	\]
	
	Using $\mathbb{E}[w_k]=0$ and
	$\mathbb{E}[w_kw_k^\top]=W$:
	\[
	\begin{aligned}
		V_k(x)=\min_u\big\{&
		x^\top Qx+u^\top Ru\\
		&+(Ax+Bu)^\top P_{k+1}(Ax+Bu)\\
		&+\operatorname{tr}(P_{k+1}W)+c_{k+1}
		\big\}.
	\end{aligned}
	\]
	
	
	Setting the derivative with respect to $u$ to zero gives
	\[
	u_k=-(R+B^\top P_{k+1}B)^{-1}
	B^\top P_{k+1}Ax_k.
	\]

		\end{flushleft}
\end{frame}

\begin{frame}{Appendix: Riccati recursion}
	\begin{flushleft}
	
	Substituting the optimal control law back into Bellman's equation,
	we obtain:
	\[
	\begin{aligned}
		P_k={}&Q+A^\top P_{k+1}A\\
		&-A^\top P_{k+1}B
		(R+B^\top P_{k+1}B)^{-1}
		B^\top P_{k+1}A,
	\end{aligned}
	\]
	
	and
	\[
	c_k=c_{k+1}+\operatorname{tr}(P_{k+1}W).
	\]
	
	The terminal conditions are
	\[
	P_N=Q_N,\qquad c_N=0.
	\]
	
	\end{flushleft}
\end{frame}





\end{document}
